\begin{frame}{Testcase 限制}
  \vspace{-0.3em}
  \begin{columns}[T,onlytextwidth]
    \column{0.48\textwidth}
    \footnotesize
    \textbf{加速基準}
    \begin{tabular}{@{}lrr@{}}
      \toprule
      Timeout & 解出 & 未解 \\
      \midrule
      30s & \BaseSolvThirty & \BaseUnsolvThirty \\
      500s & \BaseSolvFiveHundred & \BaseUnsolvFiveHundred \\
      \bottomrule
    \end{tabular}
    \column{0.50\textwidth}
    \footnotesize
    \textbf{勿混用}
    \begin{itemize}
      \item \#Ext@30s 有鑑別力
      \item \#Ext@500s $\leq$ \BaseUnsolvFiveHundred{}
      \item FALSE 可能是 timeout 標籤
      \item semantic\_extension $\neq$ \#Ext@T
      \item 一次 sample 只插一個 loop
    \end{itemize}
  \end{columns}
  \note{
    \textbf{約 45 秒}

    讀數字前要先理解 benchmark：在 500 秒內 baseline 只剩 3 題沒解，所以 \#Ext@500s 上限就是 3，鑑別力很低。真正有訊息的是 30 秒：有 118 題 baseline 沒解。另外 baseline 的 FALSE 有時是 timeout 分類，不是反例；還有 11 題 semantic extension，和官方 \#Ext 不是同一回事。還要記得 pipeline 限制：每個 sample 只在模型自選的一個 loop 插一條 assume，多 loop 題不保證每個 loop 都被抽到。
  }
\end{frame}
